\documentclass[oneside,a4paper,14pt]{extarticle}
\usepackage[a4paper,top=20mm,bottom=20mm,left=30mm,right=10mm]{geometry}
\usepackage{amsmath, amssymb, amsfonts, amsthm}
\usepackage[T2A]{fontenc}
\usepackage[utf8]{inputenc}
\usepackage[russian]{babel}
\usepackage{textcomp}
\usepackage{indentfirst}
\usepackage{graphicx}
\usepackage{float}
\usepackage{wrapfig}
\usepackage{caption}
\usepackage{tocloft}
\usepackage{seqsplit}
\allowbreak
\renewcommand{\cftsecleader}{\cftdotfill{\cftdotsep}}
\usepackage[breaklinks]{hyperref}
\usepackage{titlesec}
\usepackage{fancyvrb}
\usepackage{multirow}
\usepackage{xcolor}
\titleformat{\section}{\normalsize\bfseries}{\thesection}{1em}{}
\titleformat{\section}{\normalsize\bfseries}{\thesection}{1em}{}
\titleformat{\subsection}{\normalsize\bfseries}{\thesubsection}{1em}{}
\renewcommand\baselinestretch{1.33}\normalsize
\setlength{\parindent}{1.25cm}

\captionsetup[figure]{labelsep=period}

\begin{document}
\newpage\thispagestyle{empty}
\newpage
\setcounter{page}{1}
\tableofcontents
\newpage

\section*{Задание}
\addcontentsline{toc}{section}{Задание}
Разработать микропрограмму для выполнения заданной операции в ОУМ. Микропрограмма должна выполнять сдвиг 8-разрядного двоичного кода $X$ на $N$ разрядов в сторону старших разрядов, сдвиг логический.

\section*{Разработка алгоритма}
\addcontentsline{toc}{section}{Разработка алгоритма}
Для выполнения сдвига 8-разрядного двоичного кода $X$ на $N$ разрядов в сторону старших разрядов, сдвиг логический, необходимо выполнить алгоритм, представленный на рисунке \ref{fig:alg}.
\begin{figure}[H]
	\centering
	\includegraphics[width=0.75\textwidth]{Алгоритм.png}
	\caption{Алгоритм микропрограммы}
	\label{fig:alg}
\end{figure}
B и N – вспомогательные переменные
Перед началом цикла сдвига присваивается нулевое значение признаку равенства нуля
(E);
В цикле сдвига проверяется значение N. Если N=0, то признаку равенства нуля устанавливается значение 1, и цикл завершается;
Если N>0, то выполняется сдвиг X в сторону старших разрядов;
После сдвига значение N уменьшается на 1.

\section*{Разработка граф-схемы микропрограммы}
\addcontentsline{toc}{section}{Разработка граф-схемы микропрограммы}
Граф-схема микропрограммы представлена на рисунке \ref{fig:graph}.
\begin{figure}[H]
	\centering
	\includegraphics[width=0.75\textwidth]{Граф-схема.png}
	\caption{Граф-схема микропрограммы}
	\label{fig:graph}
\end{figure}

\newpage

\section*{Микрокоманды}
\addcontentsline{toc}{section}{Микрокоманды}
Результатом выполнения лабораторной работы является последовательность микрокоманд для выполнения в имитаторе операционного устройства. Микрокоманды представлены в таблице 1.
\begin{flushright}
	Таблица 1: Таблица микрокоманд
\end{flushright}

\begin{table}[h!]
	\centering
	\renewcommand{\arraystretch}{1.2}
	\resizebox{\textwidth}{!}{
		\begin{tabular}{|c|c|c|c|c|c|c|c|c|c|c|}

			\hline

			\multirow{2}{*}{№} & \multicolumn{3}{c|}{Код микрокоманды} & \multirow{2}{*} {Шина данных} & \multirow{2}{*}{Адрес} & \multirow{2}{*}{Адрес} & \multicolumn{4}{c|}{Управляющие сигналы}                                               \\ \cline{2-4} \cline{8-11}
			                   & 18-16                                 & 15-13                         & 12-10                  &                        & A                                        & B    & C0 & \textasciicircum OE & SC1 & SC0 \\ \hline
			00:                & 011                                   & 011                           & 111                    & 10000001               & 0000                                     & 0000 & 0  & 1                   & 0   & 0   \\ \hline
			01:                & 011                                   & 011                           & 111                    & 00000001               & 0000                                     & 1001 & 0  & 1                   & 0   & 0   \\ \hline
			02:                & 111                                   & 011                           & 000                    & 00000000               & 0000                                     & 0000 & 0  & 1                   & 0   & 0   \\ \hline
			03:                & 011                                   & 001                           & 011                    & 00000000               & 0000                                     & 1001 & 0  & 1                   & 0   & 0   \\ \hline
		\end{tabular}
	}
\end{table}
\newpage

\section*{Экранные формы}
\addcontentsline{toc}{section}{Экранные формы}
На рисунках \ref{fig:screen1} и \ref{fig:screen2} представлены экранные формы имитатора операционного устройства,в который была введена разработанная микропрограмма. На рисунке \ref{fig:screen1} представлено окно выполнения микрокоманд, а на рисунке \ref{fig:screen2} представлено окно ввода микрокоманд.
\begin{figure}[H]
	\centering
	\includegraphics[width=0.75\textwidth]{Окно выполнения микрокоманд.png}
	\caption{Окно выполнения микрокоманд}
	\label{fig:screen1}
\end{figure}

\begin{figure}[H]
	\centering
	\includegraphics[width=0.75\textwidth]{Окно ввода микрокоманд.png}
	\caption{Окно ввода микрокоманд}
	\label{fig:screen2}
\end{figure}

\section*{Вывод}
\addcontentsline{toc}{section}{Вывод}
В данной лабораторной работе была разработана микропрограмма для выполнения
сдвига 8-разрядного двоичного кода X на N разрядов в сторону старших разрядов (сдвиг логический). Были разработаны алгоритм, граф-схема и экранные формы микропрограммы.

\end{document}